% problem-000468 9 Wittig类反应拓展
\section*{9. Wittig类反应拓展}
Wittig 于 1953 年报道了磷叶⽴德（Phosphoniumylide，⼜称 Wittig 试剂）与醛、酮反应，将羰基直接转化为碳碳双键，同时⽣成副产物三苯氧膦，称之为 Wittig 反应，是构建碳碳双键的重要⽅法。有关改进和拓展Wittig反应的研究得到了⼴泛的重视和发展，⽂献报道了如下所⽰的类似Wittig反应的研究：

（图：）

\subsection*{9-1}
给出反应物中C=N双键的构型（�或�）。

\subsection*{9-2}
画出反应物 $\mathrm { P h _ { 3 } P { = } C H C N }$ 的共振结构式。

\subsection*{9-3}
指出下列哪⼀种1HNMR谱信息可⽤来测定反应物C=N双键构型。I. 化学位移 II. 磁各向异性 III. 屏蔽效应 IV. NOE 效应

\subsection*{9-4}
说明反应物结构中的- $- S O _ { 2 } \mathrm { M e }$ 基团对该反应的作⽤。

\subsection*{9-5}
研究上述反应机理时，发现如下反应：

（图：）

分别给出⽣成 ${ { \bf { M } } _ { 2 } }$ 和 ${ { \bf { M } } _ { 3 } }$ 的反应机理。

\subsection*{9-6}
在9-5反应中，以四氢呋喃(THF)为溶剂，室温下反应后，再⽤福尔⻢林(formalin)淬灭反应，得到如下实验结果：

（图：）

写出 $\mathbf { P } _ { 1 \mathrm { s } }$ ${ \bf P } _ { 2 }$ 的结构简式。

\subsection*{9-7}
Wittig 反应⼀般是在⾮质⼦极性溶剂中进⾏，但在 9-6 反应中，⽤福尔⻢林(formalin)淬灭反应得到了产物 ${ \bf P } _ { 1 }$ 和 $\mathbf { P } _ { 2 }$ ，写出 $\mathbf { P } _ { 2 }$ 形成过程中关键中间体的结构简式。
